\section{Conclusion}
RePressDB coordinates operation-dependent query distributions with empirical tail admission and recoverable physical replacement. Consecutive feasible decision epochs connect historical evidence to the current risk criterion, while directed installation-plus-return prices account for a later repair action. Exact conditional counts reduce online decision work, and three durable catalog transactions preserve activation and retirement ordering.

The common-executor evaluation exposes a specific trade-off: migrations fall from 155 to 66 and event p99 relative to cumulative-credit admission has ratio 0.872 (0.816--0.928), while the service-time ratio is 1.029 (0.995--1.065). Persistent shifts yield a 25.0\% service reduction against Static; continuous mixtures retain their observed event-tail cost. Decision records, post-action model opportunity, complete replacement measurements, and 492 successful software-fault tests connect these results to the executed mechanism. Calibration reuse and the measured data/cache conditions define the lifecycle in which these operating choices can be assessed.
